\documentclass[11pt]{article}
\usepackage{preamble}
\usepackage{amssymb}

\newcommand{\lessonrow}[2]{\textbf{#1} & #2 \\ \hline}
\newcommand{\logrulescard}{%
\textbf{PROPERTIES OF LOGARITHMS (all for \(b > 0\), \(b \neq 1\); \(M, N > 0\)):}\\
\hspace*{1.4em}PRODUCT: \(\log_b(MN) = \log_b M + \log_b N\)\\
\hspace*{1.4em}QUOTIENT: \(\log_b(M/N) = \log_b M - \log_b N\)\\
\hspace*{1.4em}POWER: \(\log_b(M^n) = n \cdot \log_b M\)\\
\hspace*{1.4em}IDENTITY: \(\log_b b = 1\) \hspace{1.2em} \(\log_b 1 = 0\) \hspace{1.2em} \(b^{\log_b x} = x\)\\
\textbf{CHANGE OF BASE:}\\
\hspace*{1.4em}\(\log_b x = (\log x) / (\log b) = (\ln x) / (\ln b)\)\\
\textbf{SOLVING EXPONENTIAL EQUATIONS:}\\
\hspace*{1.4em}\(b^x = k \Rightarrow x = \log_b k = (\log k) / (\log b)\)%
}
\newcommand{\richterprompt}{%
Performance Task The magnitude, or intensity, of an earthquake is measured on the Richter scale. For an earthquake where the amplitude of its seismographic trace is \(A\), its magnitude is modeled by the function:
\[
R(A)=\log(A)/(A_0),
\]
where \(A_0\) represents the amplitude of the smallest detectable earthquake.

\textbf{Part A} An earthquake occurs with an amplitude 200 times greater than the amplitude of the smallest detectable earthquake, \(A_0\). What is the magnitude of this earthquake on the Richter scale?

\textbf{Part B} Approximately how many times as great is the amplitude of an earthquake measuring 6.8 on the Richter scale than the amplitude of an earthquake measuring 5.9 on the Richter scale?

\textbf{Part C} Suppose the intensity of one earthquake is 150 times as great as that of another. How much greater is the magnitude of the more intense earthquake than the less intense earthquake?%
}

\begin{document}

\packetheading{Lesson 6-5 \textperiodcentered\ Quick \textperiodcentered\ Student Packet}{Properties of Logarithms + DOK-3 Power Property Proof \& Richter Generalization}

\sectionbanner{OBJECTIVES}
{\small
\renewcommand{\arraystretch}{1.25}
\noindent\begin{tabular}{|p{1.8in}|p{4.8in}|}
\hline
\lessonrow{Math Objective}{Use the Product, Quotient, and Power Properties of Logarithms, plus the Change of Base Formula, to expand, condense, and evaluate logarithmic expressions and solve exponential equations.}
\lessonrow{Language Objective}{State each log property in `words' form (`the log of a product equals the sum of the logs'), and justify a condensing step by naming which property applies.}
\lessonrow{Essential Understanding}{The log properties come directly from the exponent laws --- every property of logarithms IS a property of exponents written in inverse form.}
\end{tabular}
}

\sectionbanner{TOPIC GOALS / ESSENTIAL QUESTION / MATERIALS}
{\small
\renewcommand{\arraystretch}{1.25}
\noindent\begin{tabular}{|p{1.8in}|p{4.8in}|}
\hline
\lessonrow{Topic Goals}{Fluency with all four log properties + change of base; apply to condense, expand, and solve; prove Power Property from first principles (DOK-3).}
\lessonrow{Essential Question}{How do the properties of exponents become the properties of logarithms, and how do we use them to simplify or solve equations?}
\lessonrow{Materials}{Student packet, pencil, calculator. DOK-3 paired driver (\#13 proof + \#40 Richter) is the lesson capstone and rehearses Form B DOK-3 generalization reasoning. Note: Form B Q8 and Q15 (geometric series) are NOT on the Topic 6 assessment and are NOT covered here.}
\end{tabular}
}

\sectionbanner{LAUNCH --- Example 3: Write as a Single Logarithm (teacher models)}

\begin{callout}{\IconBook\ PROPERTIES OF LOGARITHMS (keep printed on your page)}{calloutgreen}
\logrulescard
\end{callout}

\begin{callout}{\IconWarn\ COMPRESSED LESSON --- ONE EXAMPLE IN LAUNCH, ALL THREE PROPERTIES ACTIVE}{calloutred}
Teacher models Example 3 (condensing into a single logarithm). Condensing forces you to recognize Product, Quotient, AND Power Properties at once --- watch which property applies at each step.

The DOK-3 task (\#13) asks you to prove the Power Property from scratch using exponent laws. That proof IS on Form B. Plan 12 min for \#13 + \#40.
\end{callout}

\sectionbanner{EXPLORE --- Think-Pair-Share (28 min)}
\textit{Work with a partner. Move in order. \#13 + \#40 are the big ones --- plan \(\sim\)12 min for the paired DOK-3 driver. On Wed-F 45-min, skip \#3.}

\textbf{Bank-item labels (in order):}
\begin{enumerate}[leftmargin=1.6em,itemsep=2pt,topsep=2pt]
  \item Practice \#19 --- Condense: \(\log_5 6 + \dfrac{1}{2}\log_5 y\)
  \item Practice \#21 --- Condense: \(\dfrac{1}{3}\ln 27 - 3\ln(2y)\)
  \item Practice \#34 --- Solve: \(7^x = 100\) (change of base)
  \item Practice \#3 --- Amanda's expand error on \(\log_4(c^2 d^5)\) [SKIP on Wed-F 45-min]
  \item Practice \#13 \IconStar\ DOK-3 PROOF --- Prove the Power Property: \(\log_b(M^n) = n \cdot \log_b M\)
  \item Practice \#40 Part C \IconStar\ DOK-3 RICHTER --- Generalize: one earthquake \(150\times\) greater; how much larger is its Richter magnitude?
\end{enumerate}

\bankitem{Practice \#19 --- Condense: \(\log_5 6 + \dfrac{1}{2}\log_5 y\)}{%
Use the properties of logarithms to write each expression as a single logarithm.
\[
\log_5 6+\dfrac{1}{2}\log_5 y
\]
}{}

\bankitem{Practice \#21 --- Condense: \(\dfrac{1}{3}\ln 27 - 3\ln(2y)\)}{%
Use the properties of logarithms to write each expression as a single logarithm.
\[
\dfrac{1}{3}\ln27-3\ln(2y)
\]
}{}

\bankitem{Practice \#34 --- Solve: \(7^x = 100\) (change of base)}{%
Use the Change of Base Formula to solve each equation for \(x\). Give an exact solution as a logarithm and an approximate solution rounded to the nearest thousandth.
\[
7^x=100
\]
}{}

\bankitem{Practice \#3 --- Amanda's expand error on \(\log_4(c^2 d^5)\) [SKIP on Wed-F 45-min]}{%
Error Analysis Amanda used the properties of logarithms to expand \(\log_4(c^2 d^5)\). Her work is shown below. Describe and correct the error.
\[
\log_4(c^2 d^5) = 2 \cdot \log_4 c \cdot 5 \cdot \log_4 d
\]
}{}

\bankitem{Practice \#13 \IconStar\ DOK-3 PROOF --- Prove the Power Property: \(\log_b(M^n) = n \cdot \log_b M\)}{%
Use the properties of exponents to prove the Power Property of Logarithms.
}{}

\bankitem{Practice \#40 Part C \IconStar\ DOK-3 RICHTER --- Generalize: one earthquake \(150\times\) greater; how much larger is its Richter magnitude?}{%
\richterprompt
}{}

\sectionbanner{SHARE / SUMMARY (5 min)}

\begin{sentenceframebox}
\textbf{Sentence frames}

Pair share (Power Property proof): ``Let \(x = \log_b M\). Then \(b^x = M\). Raising both sides to the nth power: \((b^x)^n = M^n\), so \(b^{nx} = M^n\). By the definition of log, \(\log_b(M^n) = nx = n \cdot \log_b M\). \(\checkmark\)''

Richter: ``\(R(150A_0) - R(A_0) = \log(150A_0/S) - \log(A_0/S) = \log 150 \approx 2.18\). So the larger earthquake is about 2.18 Richter units greater.''

Whole class: one pair reads the proof. Class signals thumbs up/sideways.
\end{sentenceframebox}

\begin{callout}{\IconSearch\ BRIDGE TO TOPIC 6 ASSESSMENT}{calloutpeach}
The Power Property proof (q13) and the Richter generalization (q40 Part C) are the exact DOK-3 reasoning pattern on Form B. You have now rehearsed both. Note: Form B Q8 and Q15 (geometric series) are NOT on the Topic 6 assessment.
\end{callout}

\sectionbanner{EXIT}

\begin{summaryexitbox}{EXIT TICKET --- Two-Part Summary}
A biologist models a population as \(P = 3 \log t + \log 4\).

(a) Write \(P\) as a single logarithm: \(P = \log(\text{\_\_\_})\)

(b) Solve for \(t\) when \(P = 2\): \(t = \text{\_\_\_}\)

Expected: (a) \(P = \log(4t^3)\) \hspace{1.8em} (b) \(t = (25)^{1/3} \approx 2.924\)

One biggest thing I learned today: \_\_\_\_\_\_\_\_\_\_\_\_\_\_\_\_\_\_\_\_\_\_\_\_\_\_\_\_\_\_
\end{summaryexitbox}

\clearpage

\sectionbanner{OPTIONAL REINFORCEMENT (not collected)}
\textit{If you finish \#13 + \#40 early. Part A of \#40 is a direct numerical Richter computation --- builds toward the Part C generalization.}

\bankitem{Optional  (Savvas Practice \#40 Part A --- direct Richter computation)}{%
\richterprompt
}{}

\end{document}
